\clearpage

\section{Fitter performance}
\label{sec:appendix:fitter_performance}

The performance of the fitter has been evaluated at various oscillation points, using various POT, various levels of throwing systematics, and various fitters.
For each study, 10000 toy datasets were created, and were fitted using MINUIT.
Residual and pull distributions are presented in the following sections.
The following three pull definitions are used:
\begin{equation}
\textrm{Pull}_{\textrm{\footnotesize{1}}} = \frac{\textrm{(best-fit parameter) - (input parameter)}}{\sigma_{\textrm{\small{fit}}}}
\end{equation}
\begin{equation}
\textrm{Pull}_{\textrm{\footnotesize{2}}} = \frac{\textrm{(best-fit parameter) - (input parameter)}}
{\sqrt{\sigma_{\textrm{\small{prior}}}^{2} - \sigma_{\textrm{\small{fit}}}^{2}}},
\end{equation}
and
\begin{equation}
\textrm{Pull}_{\textrm{\footnotesize{3}}} = \frac{\textrm{(best-fit parameter) - (nominal parameter)}}{\sigma_{\textrm{\small{fit}}}}
\end{equation}
where 
$\sigma_{\textrm{\small{fit}}}$ is output MINUIT/HESSE error and
$\sigma_{\textrm{\small{prior}}}$ is the prior assumed error.
When $\textrm{Pull}_{\textrm{\footnotesize{2}}}$ is shown for a correlated systematic parameter, 
$\sigma_{\textrm{\small{prior}}}$ is defined as the square root of the corresponding diagonal element of the covariance matrix.

The various pull definitions are used as shown below

\begin{itemize}
\item
Definition 1 is used for the pulls of both oscillation and systematic parameters in toy-MC studies.
\item
Definition 2 is used for the pulls of systematic parameters in toy-MC studies.
\item
Definition 3 is used for the pulls of systematic parameters both in toy-MC studies and for the fit to the actual data.
\end{itemize}

\clearpage
\subsection{'2+0' fit, no systematic fluctations}

Figs.\ \ref{fig:pulls:20:nosyst:mix23pull}-\ref{fig:pulls:20:nosyst:mix23} show the pull and fitted values of
$\sin^22\theta_{23}$.
Figs.\ \ref{fig:pulls:20:nosyst:dm23pull}-\ref{fig:pulls:20:nosyst:dm23} show the corresponding plots for 
$|\Delta m^2_{32}|$.
The input values are $\sin^22\theta_{23}$ = 0.9, $|\Delta m^2_{32}|$ = $2.4 \times 10^{-3}$ $eV^2/c^4$.
A bias towards higher $\sin^22\theta_{23}$ can be seen, which is present even when the pileup near the physical boundary is removed.
This bias is due to low statistics at 3.01$\times$20 POT. When Poisson fluctuations are made of a small number of events in a bin of 
the reconstructed energy distribution, a downward fluctuation is more likely than an upward fluctuation. For example, a Poisson 
distribution with mean 3 has P(1) = 0.149 and P(2) = 0.224, whereas P(4) = 0.168 and P(5) = 0.101. A downward fluctuation in a few 
bins around the oscillation maximum increases the fitted value of $\sin^{2}(2\theta_{23})$, and this is the reason for the bias 
towards high values of this parameter.

\begin{figure*}[htb]
\begin{minipage}[t]{0.48\linewidth}
\center
\includegraphics[width=20pc]{{./images/pulls/20_non_0.900000_mixing23_pull}.eps}
\caption {
  $\textrm{Pull}_{\textrm{\footnotesize{1}}}$ of sin$^22\theta_{23}$ for a 2+0 fit, with oscillation parameters
  $\sin^22\theta_{12} = $0.857; $\sin^22\theta_{13} = $0.098; $\sin^22\theta_{23} = $0.9; $\delta_{CP} = $        0; $\Delta m^2_{21} = $    7.5e-05 $eV^2/c^4$; $|\Delta m^2_{32}| = $    0.0024 $eV^2/c^4$
  and for an exposure of 3.010$\times$10$^{20}$ POT. No systematic fluctuations are included.
}
\label{fig:pulls:20:nosyst:mix23pull}
\end{minipage}
\hspace{0.06\linewidth}
\begin{minipage}[t]{0.48\linewidth}
\center
\includegraphics[width=20pc]{{./images/pulls/20_non_0.900000_mixing23}.eps}
\caption {
  Fitted value of sin$^22\theta_{23}$ for a 2+0 fit, with oscillation parameters
  $\sin^22\theta_{12} = $0.857; $\sin^22\theta_{13} = $0.098; $\sin^22\theta_{23} = $0.9; $\delta_{CP} = $        0; $\Delta m^2_{21} = $    7.5e-05 $eV^2/c^4$; $|\Delta m^2_{32}| = $    0.0024 $eV^2/c^4$
  and for an exposure of 3.010$\times$10$^{20}$ POT. No systematic fluctuations are included.
}
\label{fig:pulls:20:nosyst:mix23}
\end{minipage}
\vfill
\begin{minipage}[t]{0.48\linewidth}
\center
\includegraphics[width=20pc]{{./images/pulls/20_non_0.900000_dm23sq_pull}.eps}
\caption {
  $\textrm{Pull}_{\textrm{\footnotesize{1}}}$ of $|\Delta m^2_{32}|$ for a 2+0 fit, with oscillation parameters
  $\sin^22\theta_{12} = $0.857; $\sin^22\theta_{13} = $0.098; $\sin^22\theta_{23} = $0.9; $\delta_{CP} = $        0; $\Delta m^2_{21} = $    7.5e-05 $eV^2/c^4$; $|\Delta m^2_{32}| = $    0.0024 $eV^2/c^4$
  and for an exposure of 3.010$\times$10$^{20}$ POT. No systematic fluctuations are included.
}
\label{fig:pulls:20:nosyst:dm23pull}
\end{minipage}
\hspace{0.06\linewidth}
\begin{minipage}[t]{0.48\linewidth}
\center
\includegraphics[width=20pc]{{./images/pulls/20_non_0.900000_dm23sq}.eps}
\caption {
  Fitted value of $|\Delta m^2_{32}|$ for a 2+0 fit, with oscillation parameters
  $\sin^22\theta_{12} = $0.857; $\sin^22\theta_{13} = $0.098; $\sin^22\theta_{23} = $0.9; $\delta_{CP} = $        0; $\Delta m^2_{21} = $    7.5e-05 $eV^2/c^4$; $|\Delta m^2_{32}| = $    0.0024 $eV^2/c^4$
  and for an exposure of 3.010$\times$10$^{20}$ POT. No systematic fluctuations are included.
}
\label{fig:pulls:20:nosyst:dm23}
\end{minipage}
\end{figure*}


% \clearpage
% \subsection{'2+0' fit, systematic fluctations}

% Figs.\ \ref{fig:pulls:20:allsyst:mix23pull}-\ref{fig:pulls:20:allsyst:mix23} show the pull and fitted values of
% $\sin^22\theta_{23}$.
% Figs.\ \ref{fig:pulls:20:allsyst:dm23pull}-\ref{fig:pulls:20:allsyst:dm23} show the corresponding plots for 
% $\Delta m^2_{32}$.
% The input values are $\sin^22\theta_{23}$ = 0.9, $\Delta m^2_{32}$ = $2.4 \times 10^{-3}$.
% A bias towards higher $\sin^22\theta_{23}$ can be seen, which is present even when the pileup near the physical boundary is removed.

% \begin{figure*}[htb]
% \begin{minipage}[t]{0.48\linewidth}
% \center
% \includegraphics[width=20pc]{{./images/pulls/20_all_0.900000_mixing23_pull}.eps}
% \caption {
%   \textrm{Pull}_{\textrm{\footnotesize{1}}} of sin$^22\theta_{23}$ for a 2+0 fit, with oscillation parameters
%   $\sin^22\theta_{12} = $0.857; $\sin^22\theta_{13} = $0.098; $\sin^22\theta_{23} = $0.9; $\delta_{CP} = $        0; $\Delta m^2_{12} = $    7.5e-05; $\Delta m^2_{32} = $    0.0024
%   and for an exposure of 3.010$\times$10$^{20}$ POT. All systematic fluctuations are included.
% }
% \label{fig:pulls:20:allsyst:mix23pull}
% \end{minipage}
% \hspace{0.06\linewidth}
% \begin{minipage}[t]{0.48\linewidth}
% \center
% \includegraphics[width=20pc]{{./images/pulls/20_all_0.900000_mixing23}.eps}
% \caption {
%   Fitted value of sin$^22\theta_{23}$ for a 2+0 fit, with oscillation parameters
%   $\sin^22\theta_{12} = $0.857; $\sin^22\theta_{13} = $0.098; $\sin^22\theta_{23} = $0.9; $\delta_{CP} = $        0; $\Delta m^2_{12} = $    7.5e-05; $\Delta m^2_{32} = $    0.0024
%   and for an exposure of 3.010$\times$10$^{20}$ POT. All systematic fluctuations are included.
% }
% \label{fig:pulls:20:allsyst:mix23}
% \end{minipage}
% \vfill
% \begin{minipage}[t]{0.48\linewidth}
% \center
% \includegraphics[width=20pc]{{./images/pulls/20_all_0.900000_dm23sq_pull}.eps}
% \caption {
%   \textrm{Pull}_{\textrm{\footnotesize{1}}} of $\Delta m^2_{32}$ for a 2+0 fit, with oscillation parameters
%   $\sin^22\theta_{12} = $0.857; $\sin^22\theta_{13} = $0.098; $\sin^22\theta_{23} = $0.9; $\delta_{CP} = $        0; $\Delta m^2_{12} = $    7.5e-05; $\Delta m^2_{32} = $    0.0024
%   and for an exposure of 3.010$\times$10$^{20}$ POT. All systematic fluctuations are included.
% }
% \label{fig:pulls:20:allsyst:dm23pull}
% \end{minipage}
% \hspace{0.06\linewidth}
% \begin{minipage}[t]{0.48\linewidth}
% \center
% \includegraphics[width=20pc]{{./images/pulls/20_all_0.900000_dm23sq}.eps}
% \caption {
%   Fitted value of $\Delta m^2_{32}$ for a 2+0 fit, with oscillation parameters
%   $\sin^22\theta_{12} = $0.857; $\sin^22\theta_{13} = $0.098; $\sin^22\theta_{23} = $0.9; $\delta_{CP} = $        0; $\Delta m^2_{12} = $    7.5e-05; $\Delta m^2_{32} = $    0.0024
%   and for an exposure of 3.010$\times$10$^{20}$ POT. All systematic fluctuations are included.
% }
% \label{fig:pulls:20:allsyst:dm23}
% \end{minipage}
% \end{figure*}


\clearpage
\subsection{'2+0' fit, no systematic fluctations, 100$\times$ RunI-III POT}

Figs.\ \ref{fig:pulls:20:highPOT:nosyst:mix23pull}-\ref{fig:pulls:20:highPOT:nosyst:mix23} show the pull and fitted values of
$\sin^22\theta_{23}$.
Figs.\ \ref{fig:pulls:20:highPOT:nosyst:dm23pull}-\ref{fig:pulls:20:highPOT:nosyst:dm23} show the corresponding plots for 
$|\Delta m^2_{32}|$.
The input values are $\sin^22\theta_{23}$ = 0.9, $|\Delta m^2_{32}|$ = $2.4 \times 10^{-3}$ $eV^2/c^4$.
The bias in $\sin^22\theta_{23}$ seen at current POT has gone.

\begin{figure*}[htb]
\begin{minipage}[t]{0.48\linewidth}
\center
\includegraphics[width=20pc]{{./images/pulls/20_non_highpot_0.900000_mixing23_pull}.eps}
\caption {
  $\textrm{Pull}_{\textrm{\footnotesize{1}}}$ of sin$^22\theta_{23}$ for a 2+0 fit, with oscillation parameters
  $\sin^22\theta_{12} = $0.857; $\sin^22\theta_{13} = $0.098; $\sin^22\theta_{23} = $0.9; $\delta_{CP} = $        0; $\Delta m^2_{21} = $    7.5e-05 $eV^2/c^4$; $|\Delta m^2_{32}| = $    0.0024 $eV^2/c^4$
  and for an exposure of 3.010$\times$10$^{22}$ POT. No systematic fluctuations are included.
}
\label{fig:pulls:20:highPOT:nosyst:mix23pull}
\end{minipage}
\hspace{0.06\linewidth}
\begin{minipage}[t]{0.48\linewidth}
\center
\includegraphics[width=20pc]{{./images/pulls/20_non_highpot_0.900000_mixing23}.eps}
\caption {
  Fitted value of sin$^22\theta_{23}$ for a 2+0 fit, with oscillation parameters
  $\sin^22\theta_{12} = $0.857; $\sin^22\theta_{13} = $0.098; $\sin^22\theta_{23} = $0.9; $\delta_{CP} = $        0; $\Delta m^2_{21} = $    7.5e-05 $eV^2/c^4$; $|\Delta m^2_{32}| = $    0.0024 $eV^2/c^4$
  and for an exposure of 3.010$\times$10$^{22}$ POT. No systematic fluctuations are included.
}
\label{fig:pulls:20:highPOT:nosyst:mix23}
\end{minipage}
\vfill
\begin{minipage}[t]{0.48\linewidth}
\center
\includegraphics[width=20pc]{{./images/pulls/20_non_highpot_0.900000_dm23sq_pull}.eps}
\caption {
  $\textrm{Pull}_{\textrm{\footnotesize{1}}}$ of $|\Delta m^2_{32}|$ for a 2+0 fit, with oscillation parameters
  $\sin^22\theta_{12} = $0.857; $\sin^22\theta_{13} = $0.098; $\sin^22\theta_{23} = $0.9; $\delta_{CP} = $        0; $\Delta m^2_{21} = $    7.5e-05 $eV^2/c^4$; $|\Delta m^2_{32}| = $    0.0024 $eV^2/c^4$
  and for an exposure of 3.010$\times$10$^{22}$ POT. No systematic fluctuations are included.
}
\label{fig:pulls:20:highPOT:nosyst:dm23pull}
\end{minipage}
\hspace{0.06\linewidth}
\begin{minipage}[t]{0.48\linewidth}
\center
\includegraphics[width=20pc]{{./images/pulls/20_non_highpot_0.900000_dm23sq}.eps}
\caption {
  Fitted value of $|\Delta m^2_{32}|$ for a 2+0 fit, with oscillation parameters
  $\sin^22\theta_{12} = $0.857; $\sin^22\theta_{13} = $0.098; $\sin^22\theta_{23} = $0.9; $\delta_{CP} = $        0; $\Delta m^2_{21} = $    7.5e-05 $eV^2/c^4$; $|\Delta m^2_{32}| = $    0.0024 $eV^2/c^4$
  and for an exposure of 3.010$\times$10$^{22}$ POT. No systematic fluctuations are included.
}
\label{fig:pulls:20:highPOT:nosyst:dm23}
\end{minipage}
\end{figure*}


\clearpage
\subsection{'2+8' fit, no systematic fluctations}

Figs.\ \ref{fig:pulls:28:nosyst:mix23pull}-\ref{fig:pulls:28:nosyst:mix23} show the pull and fitted values of
$\sin^22\theta_{23}$.
Figs.\ \ref{fig:pulls:28:nosyst:dm23pull}-\ref{fig:pulls:28:nosyst:dm23} show the corresponding plots for 
$|\Delta m^2_{32}|$.
The input values are $\sin^22\theta_{23}$ = 0.9, $|\Delta m^2_{32}|$ = $2.4 \times 10^{-3}$ $eV^2/c^4$.
The results are similar to the '2+0' case
(a bias towards higher $\sin^22\theta_{23}$ can be seen, which is present even when the pileup near the physical boundary is removed)
but no additional biases are present.
Therefore, adding systematics to the fit does not cause significant changes to the performance of the
oscillation part of the fit.

\begin{figure*}[htb]
\begin{minipage}[t]{0.48\linewidth}
\center
\includegraphics[width=20pc]{{./images/pulls/28_non_0.900000_mixing23_pull}.eps}
\caption {
  $\textrm{Pull}_{\textrm{\footnotesize{1}}}$ of sin$^22\theta_{23}$ for a 2+8 fit, with oscillation parameters
  $\sin^22\theta_{12} = $0.857; $\sin^22\theta_{13} = $0.098; $\sin^22\theta_{23} = $0.9; $\delta_{CP} = $        0; $\Delta m^2_{21} = $    7.5e-05 $eV^2/c^4$; $|\Delta m^2_{32}| = $    0.0024 $eV^2/c^4$
  and for an exposure of 3.010$\times$10$^{20}$ POT. No systematic fluctuations are included.
}
\label{fig:pulls:28:nosyst:mix23pull}
\end{minipage}
\hspace{0.06\linewidth}
\begin{minipage}[t]{0.48\linewidth}
\center
\includegraphics[width=20pc]{{./images/pulls/28_non_0.900000_mixing23}.eps}
\caption {
  Fitted value of sin$^22\theta_{23}$ for a 2+8 fit, with oscillation parameters
  $\sin^22\theta_{12} = $0.857; $\sin^22\theta_{13} = $0.098; $\sin^22\theta_{23} = $0.9; $\delta_{CP} = $        0; $\Delta m^2_{21} = $    7.5e-05 $eV^2/c^4$; $|\Delta m^2_{32}| = $    0.0024 $eV^2/c^4$
  and for an exposure of 3.010$\times$10$^{20}$ POT. No systematic fluctuations are included.
}
\label{fig:pulls:28:nosyst:mix23}
\end{minipage}
\vfill
\begin{minipage}[t]{0.48\linewidth}
\center
\includegraphics[width=20pc]{{./images/pulls/28_non_0.900000_dm23sq_pull}.eps}
\caption {
  $\textrm{Pull}_{\textrm{\footnotesize{1}}}$ of $|\Delta m^2_{32}|$ for a 2+8 fit, with oscillation parameters
  $\sin^22\theta_{12} = $0.857; $\sin^22\theta_{13} = $0.098; $\sin^22\theta_{23} = $0.9; $\delta_{CP} = $        0; $\Delta m^2_{21} = $    7.5e-05 $eV^2/c^4$; $|\Delta m^2_{32}| = $    0.0024 $eV^2/c^4$
  and for an exposure of 3.010$\times$10$^{20}$ POT. No systematic fluctuations are included.
}
\label{fig:pulls:28:nosyst:dm23pull}
\end{minipage}
\hspace{0.06\linewidth}
\begin{minipage}[t]{0.48\linewidth}
\center
\includegraphics[width=20pc]{{./images/pulls/28_non_0.900000_dm23sq}.eps}
\caption {
  Fitted value of $|\Delta m^2_{32}|$ for a 2+8 fit, with oscillation parameters
  $\sin^22\theta_{12} = $0.857; $\sin^22\theta_{13} = $0.098; $\sin^22\theta_{23} = $0.9; $\delta_{CP} = $        0; $\Delta m^2_{21} = $    7.5e-05 $eV^2/c^4$; $|\Delta m^2_{32}| = $    0.0024 $eV^2/c^4$
  and for an exposure of 3.010$\times$10$^{20}$ POT. No systematic fluctuations are included.
}
\label{fig:pulls:28:nosyst:dm23}
\end{minipage}
\end{figure*}

\clearpage
\subsection{'2+8' fit, dominant systematic fluctations}

Figs.\ \ref{fig:pulls:28:selsyst:mix23pull}-\ref{fig:pulls:28:selsyst:mix23} show the pull and fitted values of
$\sin^22\theta_{23}$.
Figs.\ \ref{fig:pulls:28:selsyst:dm23pull}-\ref{fig:pulls:28:selsyst:dm23} show the corresponding plots for 
$|\Delta m^2_{32}|$.
The input values are $\sin^22\theta_{23}$ = 0.9, $|\Delta m^2_{32}|$ = $2.4 \times 10^{-3}$ $eV^2/c^4$.
The bias in $\sin^22\theta_{23}$ is still present.

Fig.\ \ref{fig:pulls:28:selsyst:sysresidual} shows the residuals of the systematic parameters.
Figs.\ \ref{fig:pulls:28:selsyst:syspull}-\ref{fig:pulls:28:selsyst:sysaltpull2} show the pulls distributions of the systematic parameters.
Some of the distributions are wide, but there are no unexpected features in the plots.
$f_{E; r}^{SK}$ in Fig. \ref{fig:pulls:28:selsyst:syspull} is non-Gaussian, but this in not unexpected;
the energy scale systematic has a discontinuous derivative, leading to issues with the Hesse error.
It should be noted that the spectral function is not expected to give a Gaussian distribution
since the throws are uniform between 0 and 1, and the fitted value is strictly confined to the range [0,1].
The plots show that a best-fit value of 0 is favoured.
%Fig.\ \ref{fig:pulls:28:selsyst:syssum} gives a summary of the mean and sigma of the fitted pulls.

\begin{figure*}[htb]
\begin{minipage}[t]{0.48\linewidth}
\center
\includegraphics[width=20pc]{{./images/pulls/28_sel_0.900000_mixing23_pull}.eps}
\caption {
  $\textrm{Pull}_{\textrm{\footnotesize{1}}}$ of sin$^22\theta_{23}$ for a 2+8 fit, with oscillation parameters
  $\sin^22\theta_{12} = $0.857; $\sin^22\theta_{13} = $0.098; $\sin^22\theta_{23} = $0.9; $\delta_{CP} = $        0; $\Delta m^2_{21} = $    7.5e-05 $eV^2/c^4$; $|\Delta m^2_{32}| = $    0.0024 $eV^2/c^4$
  and for an exposure of 3.010$\times$10$^{20}$ POT. The 8 dominant systematics are fluctuated.
}
\label{fig:pulls:28:selsyst:mix23pull}
\end{minipage}
\hspace{0.06\linewidth}
\begin{minipage}[t]{0.48\linewidth}
\center
\includegraphics[width=20pc]{{./images/pulls/28_sel_0.900000_mixing23}.eps}
\caption {
  Fitted value of sin$^22\theta_{23}$ for a 2+8 fit, with oscillation parameters
  $\sin^22\theta_{12} = $0.857; $\sin^22\theta_{13} = $0.098; $\sin^22\theta_{23} = $0.9; $\delta_{CP} = $        0; $\Delta m^2_{21} = $    7.5e-05 $eV^2/c^4$; $|\Delta m^2_{32}| = $    0.0024 $eV^2/c^4$
  and for an exposure of 3.010$\times$10$^{20}$ POT. The 8 dominant systematics are fluctuated.
}
\label{fig:pulls:28:selsyst:mix23}
\end{minipage}
\vfill
\begin{minipage}[t]{0.48\linewidth}
\center
\includegraphics[width=20pc]{{./images/pulls/28_sel_0.900000_dm23sq_pull}.eps}
\caption {
  $\textrm{Pull}_{\textrm{\footnotesize{1}}}$ of $|\Delta m^2_{32}|$ for a 2+8 fit, with oscillation parameters
  $\sin^22\theta_{12} = $0.857; $\sin^22\theta_{13} = $0.098; $\sin^22\theta_{23} = $0.9; $\delta_{CP} = $        0; $\Delta m^2_{21} = $    7.5e-05 $eV^2/c^4$; $|\Delta m^2_{32}| = $    0.0024 $eV^2/c^4$
  and for an exposure of 3.010$\times$10$^{20}$ POT. The 8 dominant systematics are fluctuated.
}
\label{fig:pulls:28:selsyst:dm23pull}
\end{minipage}
\hspace{0.06\linewidth}
\begin{minipage}[t]{0.48\linewidth}
\center
\includegraphics[width=20pc]{{./images/pulls/28_sel_0.900000_dm23sq}.eps}
\caption {
  Fitted value of $|\Delta m^2_{32}|$ for a 2+8 fit, with oscillation parameters
  $\sin^22\theta_{12} = $0.857; $\sin^22\theta_{13} = $0.098; $\sin^22\theta_{23} = $0.9; $\delta_{CP} = $        0; $\Delta m^2_{21} = $    7.5e-05 $eV^2/c^4$; $|\Delta m^2_{32}| = $    0.0024 $eV^2/c^4$
  and for an exposure of 3.010$\times$10$^{20}$ POT. The 8 dominant systematics are fluctuated.
}
\label{fig:pulls:28:selsyst:dm23}
\end{minipage}
\end{figure*}


\begin{figure*}[htb]
\center
\includegraphics[width=40pc]{{./images/pulls/28_sel_0.900000_important}.eps}
\caption {
  Residuals of the systematic parameters for a 2+8 fit, with oscillation parameters
  $\sin^22\theta_{12} = $0.857; $\sin^22\theta_{13} = $0.098; $\sin^22\theta_{23} = $0.9; $\delta_{CP} = $        0; $\Delta m^2_{21} = $    7.5e-05 $eV^2/c^4$; $|\Delta m^2_{32}| = $    0.0024 $eV^2/c^4$
  and for an exposure of 3.010$\times$10$^{20}$ POT.  The 8 dominant
  systematics are fluctuated.
  The plots correspond to (line by line, left to right):
  $f_{E; r}^{SK}$;  $f_{CCnonQE; r}^{SK}$; $f_{NC; r}^{SK}$;
  $f_{16; t; r}^{banff}$ (M$_A$QE); $f_{18; t, r}^{banff}$ (CCQE norm $<$ 1.5 GeV);
  $f_{SF; t, r}$; $f_{bindE; t, r}$; $f_{\pi-less\Delta; t, r}$.
}
\label{fig:pulls:28:selsyst:sysresidual}
\end{figure*}

\begin{figure*}[htb]
\center
\includegraphics[width=40pc]{{./images/pulls/28_sel_0.900000_important_pull}.eps}
\caption {
  $\textrm{Pull}_{\textrm{\footnotesize{2}}}$  of the systematic parameters for a 2+8 fit, with oscillation parameters
  $\sin^22\theta_{12} = $0.857; $\sin^22\theta_{13} = $0.098; $\sin^22\theta_{23} = $0.9; $\delta_{CP} = $        0; $\Delta m^2_{21} = $    7.5e-05 $eV^2/c^4$; $|\Delta m^2_{32}| = $    0.0024 $eV^2/c^4$
  and for an exposure of 3.010$\times$10$^{20}$ POT.  The 8 dominant
  systematics are fluctuated.
  The plots correspond to (line by line, left to right):
  $f_{E; r}^{SK}$;  $f_{CCnonQE; r}^{SK}$; $f_{NC; r}^{SK}$;
  $f_{16; t; r}^{banff}$ (M$_A$QE); $f_{18; t, r}^{banff}$ (CCQE norm $<$ 1.5 GeV);
  $f_{SF; t, r}$; $f_{bindE; t, r}$; $f_{\pi-less\Delta; t, r}$.
}
\label{fig:pulls:28:selsyst:syspull}
\end{figure*}

\begin{figure*}[htb]
\center
\includegraphics[width=40pc]{{./images/pulls/28_sel_0.900000_important_altpull}.eps}
\caption {
  $\textrm{Pull}_{\textrm{\footnotesize{1}}}$  of the systematic parameters for a 2+8 fit, with oscillation parameters
  $\sin^22\theta_{12} = $0.857; $\sin^22\theta_{13} = $0.098; $\sin^22\theta_{23} = $0.9; $\delta_{CP} = $        0; $\Delta m^2_{21} = $    7.5e-05 $eV^2/c^4$; $|\Delta m^2_{32}| = $    0.0024 $eV^2/c^4$
  and for an exposure of 3.010$\times$10$^{20}$ POT.  The 8 dominant
  systematics are fluctuated.
  The plots correspond to (line by line, left to right):
  $f_{E; r}^{SK}$;  $f_{CCnonQE; r}^{SK}$; $f_{NC; r}^{SK}$;
  $f_{16; t; r}^{banff}$ (M$_A$QE); $f_{18; t, r}^{banff}$ (CCQE norm $<$ 1.5 GeV);
  $f_{SF; t, r}$; $f_{bindE; t, r}$; $f_{\pi-less\Delta; t, r}$.
}
\label{fig:pulls:28:selsyst:sysaltpull}
\end{figure*}

\begin{figure*}[htb]
\center
\includegraphics[width=40pc]{{./images/pulls/28_sel_0.900000_important_altpull2}.eps}
\caption {
  $\textrm{Pull}_{\textrm{\footnotesize{3}}}$  of the systematic parameters for a 2+8 fit, with oscillation parameters
  $\sin^22\theta_{12} = $0.857; $\sin^22\theta_{13} = $0.098; $\sin^22\theta_{23} = $0.9; $\delta_{CP} = $        0; $\Delta m^2_{21} = $    7.5e-05 $eV^2/c^4$; $|\Delta m^2_{32}| = $    0.0024$eV^2/c^4$
  and for an exposure of 3.010$\times$10$^{20}$ POT.  The 8 dominant
  systematics are fluctuated.
  The plots correspond to (line by line, left to right):
  $f_{E; r}^{SK}$;  $f_{CCnonQE; r}^{SK}$; $f_{NC; r}^{SK}$;
  $f_{16; t; r}^{banff}$ (M$_A$QE); $f_{18; t, r}^{banff}$ (CCQE norm $<$ 1.5 GeV);
  $f_{SF; t, r}$; $f_{bindE; t, r}$; $f_{\pi-less\Delta; t, r}$.
}
\label{fig:pulls:28:selsyst:sysaltpull2}
\end{figure*}

% \begin{figure*}[htb]
% \center
% \includegraphics[width=30pc]{{./images/pulls/28_sel_0.900000_summary_pull}.eps}
% \caption {
%   Summary of the fitted systematics (mean and sigma of pull) for a 2+8 fit, with oscillation parameters
%   $\sin^22\theta_{12} = $0.857; $\sin^22\theta_{13} = $0.098; $\sin^22\theta_{23} = $0.9; $\delta_{CP} = $        0; $\Delta m^2_{12} = $    7.5e-05; $|\Delta m^2_{32}| = $    0.0024 $eV^2/c^4$
%   and for an exposure of 3.010$\times$10$^{20}$ POT.  The 8 dominant systematics are fluctuated.
% }
% \label{fig:pulls:28:selsyst:syssum}
% \end{figure*}


\clearpage
\subsection{'2+8' fit, all systematic fluctations}

Figs.\ \ref{fig:pulls:28:allsyst:mix23pull}-\ref{fig:pulls:28:allsyst:mix23} show the pull and fitted values of
$\sin^22\theta_{23}$.
Figs.\ \ref{fig:pulls:28:allsyst:dm23pull}-\ref{fig:pulls:28:allsyst:dm23} show the corresponding plots for 
$|\Delta m^2_{32}|$.
The input values are $\sin^22\theta_{23}$ = 0.9, $|\Delta m^2_{32}|$ = $2.4 \times 10^{-3}$ $eV^2/c^4$.
The bias in $\sin^22\theta_{23}$ is still present.

Fig.\ \ref{fig:pulls:28:allsyst:sysresidual} shows the residuals of the systematic parameters.
Figs.\ \ref{fig:pulls:28:allsyst:sys}-\ref{fig:pulls:28:allsyst:sysaltpull2} show the pull distributions of the systematic parameters.
The distributions have widened, but this is as to be expected
(the parameters that are thrown but not fitted are \emph{absorbed} by the parameters that are fitted).
There are still no unexpected features.
%Fig. \ref{fig:pulls:28:allsyst:syssum} gives a summary of the mean and sigma of the fitted pulls.


\begin{figure*}[htb]
\begin{minipage}[t]{0.48\linewidth}
\center
\includegraphics[width=20pc]{{./images/pulls/28_all_0.900000_mixing23_pull}.eps}
\caption {
  $\textrm{Pull}_{\textrm{\footnotesize{1}}}$ of sin$^22\theta_{23}$ for a 2+8 fit, with oscillation parameters
  $\sin^22\theta_{12} = $0.857; $\sin^22\theta_{13} = $0.098; $\sin^22\theta_{23} = $0.9; $\delta_{CP} = $        0; $\Delta m^2_{21} = $    7.5e-05 $eV^2/c^4$; $|\Delta m^2_{32}| = $    0.0024 $eV^2/c^4$
  and for an exposure of 3.010$\times$10$^{20}$ POT. All systematic fluctuations are included.
}
\label{fig:pulls:28:allsyst:mix23pull}
\end{minipage}
\hspace{0.06\linewidth}
\begin{minipage}[t]{0.48\linewidth}
\center
\includegraphics[width=20pc]{{./images/pulls/28_all_0.900000_mixing23}.eps}
\caption {
  Fitted value of sin$^22\theta_{23}$ for a 2+8 fit, with oscillation parameters
  $\sin^22\theta_{12} = $0.857; $\sin^22\theta_{13} = $0.098; $\sin^22\theta_{23} = $0.9; $\delta_{CP} = $        0; $\Delta m^2_{21} = $    7.5e-05 $eV^2/c^4$; $|\Delta m^2_{32}| = $    0.0024 $eV^2/c^4$
  and for an exposure of 3.010$\times$10$^{20}$ POT. All systematic fluctuations are included.
}
\label{fig:pulls:28:allsyst:mix23}
\end{minipage}
\vfill
\begin{minipage}[t]{0.48\linewidth}
\center
\includegraphics[width=20pc]{{./images/pulls/28_all_0.900000_dm23sq_pull}.eps}
\caption {
  $\textrm{Pull}_{\textrm{\footnotesize{1}}}$ of $|\Delta m^2_{32}|$ for a 2+8 fit, with oscillation parameters
  $\sin^22\theta_{12} = $0.857; $\sin^22\theta_{13} = $0.098; $\sin^22\theta_{23} = $0.9; $\delta_{CP} = $        0; $\Delta m^2_{21} = $    7.5e-05 $eV^2/c^4$; $|\Delta m^2_{32}| = $    0.0024 $eV^2/c^4$
  and for an exposure of 3.010$\times$10$^{20}$ POT. All systematic fluctuations are included.
}
\label{fig:pulls:28:allsyst:dm23pull}
\end{minipage}
\hspace{0.06\linewidth}
\begin{minipage}[t]{0.48\linewidth}
\center
\includegraphics[width=20pc]{{./images/pulls/28_all_0.900000_dm23sq}.eps}
\caption {
  Fitted value of $|\Delta m^2_{32}|$ for a 2+8 fit, with oscillation parameters
  $\sin^22\theta_{12} = $0.857; $\sin^22\theta_{13} = $0.098; $\sin^22\theta_{23} = $0.9; $\delta_{CP} = $        0; $\Delta m^2_{21} = $    7.5e-05 $eV^2/c^4$; $|\Delta m^2_{32}| = $    0.0024 $eV^2/c^4$
  and for an exposure of 3.010$\times$10$^{20}$ POT. All systematic fluctuations are included.
}
\label{fig:pulls:28:allsyst:dm23}
\end{minipage}
\end{figure*}

\begin{figure*}[htb]
\center
\includegraphics[width=40pc]{{./images/pulls/28_all_0.900000_important}.eps}
\caption {
  Residual of the systematic parameters for a 2+8 fit, with oscillation parameters
  $\sin^22\theta_{12} = $0.857; $\sin^22\theta_{13} = $0.098; $\sin^22\theta_{23} = $0.9; $\delta_{CP} = $        0; $\Delta m^2_{21} = $    7.5e-05 $eV^2/c^4$; $|\Delta m^2_{32}| = $    0.0024 $eV^2/c^4$
  and for an exposure of 3.010$\times$10$^{20}$ POT.  All systematic
  fluctuations are included.
  The plots correspond to (line by line, left to right):
  $f_{E; r}^{SK}$;  $f_{CCnonQE; r}^{SK}$; $f_{NC; r}^{SK}$;
  $f_{16; t; r}^{banff}$ (M$_A$QE); $f_{18; t, r}^{banff}$ (CCQE norm $<$ 1.5 GeV);
  $f_{SF; t, r}$; $f_{bindE; t, r}$; $f_{\pi-less\Delta; t, r}$.
}
\label{fig:pulls:28:allsyst:sysresidual}
\end{figure*}

\begin{figure*}[htb]
\center
\includegraphics[width=40pc]{{./images/pulls/28_all_0.900000_important_pull}.eps}
\caption {
  $\textrm{Pull}_{\textrm{\footnotesize{2}}}$  of the systematic parameters for a 2+8 fit, with oscillation parameters
  $\sin^22\theta_{12} = $0.857; $\sin^22\theta_{13} = $0.098; $\sin^22\theta_{23} = $0.9; $\delta_{CP} = $        0; $\Delta m^2_{21} = $    7.5e-05 $eV^2/c^4$; $|\Delta m^2_{32}| = $    0.0024 $eV^2/c^4$
  and for an exposure of 3.010$\times$10$^{20}$ POT.  All systematic
  fluctuations are included.
  The plots correspond to (line by line, left to right):
  $f_{E; r}^{SK}$;  $f_{CCnonQE; r}^{SK}$; $f_{NC; r}^{SK}$;
  $f_{16; t; r}^{banff}$ (M$_A$QE); $f_{18; t, r}^{banff}$ (CCQE norm $<$ 1.5 GeV);
  $f_{SF; t, r}$; $f_{bindE; t, r}$; $f_{\pi-less\Delta; t, r}$.
}
\label{fig:pulls:28:allsyst:sys}
\end{figure*}


\begin{figure*}[htb]
\center
\includegraphics[width=40pc]{{./images/pulls/28_all_0.900000_important_altpull}.eps}
\caption {
  $\textrm{Pull}_{\textrm{\footnotesize{1}}}$ of the systematic parameters for a 2+8 fit, with oscillation parameters
  $\sin^22\theta_{12} = $0.857; $\sin^22\theta_{13} = $0.098; $\sin^22\theta_{23} = $0.9; $\delta_{CP} = $        0; $\Delta m^2_{21} = $    7.5e-05 $eV^2/c^4$; $|\Delta m^2_{32}| = $    0.0024 $eV^2/c^4$
  and for an exposure of 3.010$\times$10$^{20}$ POT.  All systematic
  fluctuations are included.
  The plots correspond to (line by line, left to right):
  $f_{E; r}^{SK}$;  $f_{CCnonQE; r}^{SK}$; $f_{NC; r}^{SK}$;
  $f_{16; t; r}^{banff}$ (M$_A$QE); $f_{18; t, r}^{banff}$ (CCQE norm $<$ 1.5 GeV);
  $f_{SF; t, r}$; $f_{bindE; t, r}$; $f_{\pi-less\Delta; t, r}$.
}
\label{fig:pulls:28:allsyst:sysaltpull}
\end{figure*}

\begin{figure*}[htb]
\center
\includegraphics[width=40pc]{{./images/pulls/28_all_0.900000_important_altpull2}.eps}
\caption {
  $\textrm{Pull}_{\textrm{\footnotesize{3}}}$ of the systematic parameters for a 2+8 fit, with oscillation parameters
  $\sin^22\theta_{12} = $0.857; $\sin^22\theta_{13} = $0.098; $\sin^22\theta_{23} = $0.9; $\delta_{CP} = $        0; $\Delta m^2_{21} = $    7.5e-05 $eV^2/c^4$; $|\Delta m^2_{32}| = $    0.0024 $eV^2/c^4$
  and for an exposure of 3.010$\times$10$^{20}$ POT.  All systematic
  fluctuations are included.
  The plots correspond to (line by line, left to right):
  $f_{E; r}^{SK}$;  $f_{CCnonQE; r}^{SK}$; $f_{NC; r}^{SK}$;
  $f_{16; t; r}^{banff}$ (M$_A$QE); $f_{18; t, r}^{banff}$ (CCQE norm $<$ 1.5 GeV);
  $f_{SF; t, r}$; $f_{bindE; t, r}$; $f_{\pi-less\Delta; t, r}$.
}
\label{fig:pulls:28:allsyst:sysaltpull2}
\end{figure*}

% \begin{figure*}[htb]
% \center
% \includegraphics[width=30pc]{{./images/pulls/28_all_0.900000_summary_pull}.eps}
% \caption {
%   Summary of the fitted systematics (mean and sigma of pull) for a 2+8 fit, with oscillation parameters
%   $\sin^22\theta_{12} = $0.857; $\sin^22\theta_{13} = $0.098; $\sin^22\theta_{23} = $0.9; $\delta_{CP} = $        0; $\Delta m^2_{12} = $    7.5e-05; $|\Delta m^2_{32}| = $    0.0024 $eV^2/c^4$
%   and for an exposure of 3.010$\times$10$^{20}$ POT.  All systematic fluctuations are included.
% }
% \label{fig:pulls:28:allsyst:syssum}
% \end{figure*}




% \clearpage
% \subsection{'2+N' fit, no systematic fluctuations}

% Figs.\ \ref{fig:pulls:2N:nosyst:mix23pull}-\ref{fig:pulls:2N:nosyst:mix23} show the pull and fitted values of
% $\sin^22\theta_{23}$.
% Figs.\ \ref{fig:pulls:2N:nosyst:dm23pull}-\ref{fig:pulls:2N:nosyst:dm23} show the corresponding plots for 
% $\Delta m^2_{32}$.
% The input values are $\sin^22\theta_{23}$ = 0.9, $\Delta m^2_{32}$ = $2.4 \times 10^{-3}$.
% The bias in $\sin^22\theta_{23}$ is still present, but there is no significant change after including the parameters
% that are not fit for in the systematic throws.

% Fig.\ \ref{fig:pulls:2N:nosyst:sys} shows the pulls of the systematic parameters.
% The distributions have widened, but this is as to be expected
% (the parameters that are thrown but not fitted are \emph{absorbed} by the parameters that are fitted).
% There are still no strange features.
% Fig.\ \ref{fig:pulls:2N:nosyst:syssum} gives a summary of the mean and sigma of the fitted pulls.


% \begin{figure*}[htb]
% \begin{minipage}[t]{0.48\linewidth}
% \center
% \includegraphics[width=20pc]{{./images/pulls/2all_all_0.900000_mixing23_pull}.eps}
% \caption {
%   \textrm{Pull}_{\textrm{\footnotesize{1}}} of sin$^22\theta_{23}$ for a 2+N fit, with oscillation parameters
%   $\sin^22\theta_{12} = $0.857; $\sin^22\theta_{13} = $0.098; $\sin^22\theta_{23} = $0.9; $\delta_{CP} = $        0; $\Delta m^2_{12} = $    7.5e-05; $\Delta m^2_{32} = $    0.0024
%   and for an exposure of 3.010$\times$10$^{20}$ POT. No systematic fluctuations are included. 
% }
% \label{fig:pulls:2N:nosyst:mix23pull}
% \end{minipage}
% \hspace{0.06\linewidth}
% \begin{minipage}[t]{0.48\linewidth}
% \center
% \includegraphics[width=20pc]{{./images/pulls/2all_all_0.900000_mixing23}.eps}
% \caption {
%   Fitted value of sin$^22\theta_{23}$ for a 2+N fit, with oscillation parameters
%   $\sin^22\theta_{12} = $0.857; $\sin^22\theta_{13} = $0.098; $\sin^22\theta_{23} = $0.9; $\delta_{CP} = $        0; $\Delta m^2_{12} = $    7.5e-05; $\Delta m^2_{32} = $    0.0024
%   and for an exposure of 3.010$\times$10$^{20}$ POT. No systematic fluctuations are included.
% }
% \label{fig:pulls:2N:nosyst:mix23}
% \end{minipage}
% \vfill
% \begin{minipage}[t]{0.48\linewidth}
% \center
% \includegraphics[width=20pc]{{./images/pulls/2all_all_0.900000_dm23sq_pull}.eps}
% \caption {
%   \textrm{Pull}_{\textrm{\footnotesize{1}}} of $\Delta m^2_{32}$ for a 2+N fit, with oscillation parameters
%   $\sin^22\theta_{12} = $0.857; $\sin^22\theta_{13} = $0.098; $\sin^22\theta_{23} = $0.9; $\delta_{CP} = $        0; $\Delta m^2_{12} = $    7.5e-05; $\Delta m^2_{32} = $    0.0024
%   and for an exposure of 3.010$\times$10$^{20}$ POT. No systematic fluctuations are included.
% }
% \label{fig:pulls:2N:nosyst:dm23pull}
% \end{minipage}
% \hspace{0.06\linewidth}
% \begin{minipage}[t]{0.48\linewidth}
% \center
% \includegraphics[width=20pc]{{./images/pulls/2all_all_0.900000_dm23sq}.eps}
% \caption {
%   Fitted value of $\Delta m^2_{32}$ for a 2+N fit, with oscillation parameters
%   $\sin^22\theta_{12} = $0.857; $\sin^22\theta_{13} = $0.098; $\sin^22\theta_{23} = $0.9; $\delta_{CP} = $        0; $\Delta m^2_{12} = $    7.5e-05; $\Delta m^2_{32} = $    0.0024
%   and for an exposure of 3.010$\times$10$^{20}$ POT. No systematic fluctuations are included.
% }
% \label{fig:pulls:2N:nosyst:dm23}
% \end{minipage}
% \end{figure*}


% \clearpage
% \subsection{'2+N' fit, all systematic fluctuations}


% Figs.\ \ref{fig:pulls:2N:allsyst:mix23pull}-\ref{fig:pulls:2N:allsyst:mix23} show the pull and fitted values of
% $\sin^22\theta_{23}$.
% Figs.\ \ref{fig:pulls:2N:allsyst:dm23pull}-\ref{fig:pulls:2N:allsyst:dm23} show the corresponding plots for 
% $\Delta m^2_{32}$.
% The input values are $\sin^22\theta_{23}$ = 0.9, $\Delta m^2_{32}$ = $2.4 \times 10^{-3}$.
% The bias in $\sin^22\theta_{23}$ is still present, but there is no significant change after including the parameters
% that are not fit for in the systematic throws.

% Fig.\ \ref{fig:pulls:2N:allsyst:sys1}-\ref{fig:pulls:2N:allsyst:sys2} show the pulls of the systematic parameters.
% Fig.\ \ref{fig:pulls:2N:allsyst:syssum} gives a summary of the mean and sigma of the fitted pulls.


% \begin{figure*}[htb]
% \begin{minipage}[t]{0.48\linewidth}
% \center
% \includegraphics[width=20pc]{{./images/pulls/2all_all_0.900000_mixing23_pull}.eps}
% \caption {
%   \textrm{Pull}_{\textrm{\footnotesize{1}}} of sin$^22\theta_{23}$ for a 2+N fit, with oscillation parameters
%   $\sin^22\theta_{12} = $0.857; $\sin^22\theta_{13} = $0.098; $\sin^22\theta_{23} = $0.9; $\delta_{CP} = $        0; $\Delta m^2_{12} = $    7.5e-05; $\Delta m^2_{32} = $    0.0024
%   and for an exposure of 3.010$\times$10$^{20}$ POT. All systematic fluctuations are included.
% }
% \label{fig:pulls:2N:allsyst:mix23pull}
% \end{minipage}
% \hspace{0.06\linewidth}
% \begin{minipage}[t]{0.48\linewidth}
% \center
% \includegraphics[width=20pc]{{./images/pulls/2all_all_0.900000_mixing23}.eps}
% \caption {
%   Fitted value of sin$^22\theta_{23}$ for a 2+N fit, with oscillation parameters
%   $\sin^22\theta_{12} = $0.857; $\sin^22\theta_{13} = $0.098; $\sin^22\theta_{23} = $0.9; $\delta_{CP} = $        0; $\Delta m^2_{12} = $    7.5e-05; $\Delta m^2_{32} = $    0.0024
%   and for an exposure of 3.010$\times$10$^{20}$ POT. All systematic fluctuations are included.
% }
% \label{fig:pulls:2N:allsyst:mix23}
% \end{minipage}
% \vfill
% \begin{minipage}[t]{0.48\linewidth}
% \center
% \includegraphics[width=20pc]{{./images/pulls/2all_all_0.900000_dm23sq_pull}.eps}
% \caption {
%   \textrm{Pull}_{\textrm{\footnotesize{1}}} of $\Delta m^2_{32}$ for a 2+N fit, with oscillation parameters
%   $\sin^22\theta_{12} = $0.857; $\sin^22\theta_{13} = $0.098; $\sin^22\theta_{23} = $0.9; $\delta_{CP} = $        0; $\Delta m^2_{12} = $    7.5e-05; $\Delta m^2_{32} = $    0.0024
%   and for an exposure of 3.010$\times$10$^{20}$ POT. All systematic fluctuations are included.
% }
% \label{fig:pulls:2N:allsyst:dm23pull}
% \end{minipage}
% \hspace{0.06\linewidth}
% \begin{minipage}[t]{0.48\linewidth}
% \center
% \includegraphics[width=20pc]{{./images/pulls/2all_all_0.900000_dm23sq}.eps}
% \caption {
%   Fitted value of $\Delta m^2_{32}$ for a 2+N fit, with oscillation parameters
%   $\sin^22\theta_{12} = $0.857; $\sin^22\theta_{13} = $0.098; $\sin^22\theta_{23} = $0.9; $\delta_{CP} = $        0; $\Delta m^2_{12} = $    7.5e-05; $\Delta m^2_{32} = $    0.0024
%   and for an exposure of 3.010$\times$10$^{20}$ POT. All systematic fluctuations are included.
% }
% \label{fig:pulls:2N:allsyst:dm23}
% \end{minipage}
% \end{figure*}

% \begin{figure*}[htb]
% \center
% \includegraphics[width=40pc]{{./images/pulls/2all_all_0.900000_banff_pull}.eps}
% \caption {
%   Pull of the systematic parameters for a 2+N fit, with oscillation parameters
%   $\sin^22\theta_{12} = $0.857; $\sin^22\theta_{13} = $0.098; $\sin^22\theta_{23} = $0.9; $\delta_{CP} = $        0; $\Delta m^2_{12} = $    7.5e-05; $\Delta m^2_{32} = $    0.0024
%   and for an exposure of 3.010$\times$10$^{20}$ POT.  All systematic fluctuations are included.
% }
% \label{fig:pulls:2N:allsyst:sys1}
% \end{figure*}

% \begin{figure*}[htb]
% \center
% \includegraphics[width=40pc]{{./images/pulls/2all_all_0.900000_other_pull}.eps}
% \caption {
%   Pull of the systematic parameters for a 2+N fit, with oscillation parameters
%   $\sin^22\theta_{12} = $0.857; $\sin^22\theta_{13} = $0.098; $\sin^22\theta_{23} = $0.9; $\delta_{CP} = $        0; $\Delta m^2_{12} = $    7.5e-05; $\Delta m^2_{32} = $    0.0024
%   and for an exposure of 3.010$\times$10$^{20}$ POT.  All systematic fluctuations are included.
% }
% \label{fig:pulls:2N:allsyst:sys2}
% \end{figure*}

% \begin{figure*}[htb]
% \center
% \includegraphics[width=30pc]{{./images/pulls/2all_all_0.900000_summary_pull}.eps}
% \caption {
%   Summary of the fitted systematics (mean and sigma of pull) for a 2+N fit, with oscillation parameters
%   $\sin^22\theta_{12} = $0.857; $\sin^22\theta_{13} = $0.098; $\sin^22\theta_{23} = $0.9; $\delta_{CP} = $        0; $\Delta m^2_{12} = $    7.5e-05; $\Delta m^2_{32} = $    0.0024
%   and for an exposure of 3.010$\times$10$^{20}$ POT.  All systematic fluctuations are included.
% }
% \label{fig:pulls:2N:allsyst:syssum}
% \end{figure*}
